\chapter{Animação: Dando Vida aos Personagens}

\begin{objetivos}
\begin{itemize}
\item Criar animações trocando fantasias;
\item Combinar movimento e animação;
\item Fazer um personagem andar com as setas do teclado.
\end{itemize}
\end{objetivos}

\section{O segredo da animação}

Um desenho animado é feito de muitos desenhos parecidos mostrados rapidamente, um após o outro. No Scratch fazemos igual, mudando de \textbf{fantasia} em um laço, com uma pausa pequena.

\begin{pensaassim}[Pense assim: desenhos no canto do caderno]
Desenhe uma pessoa em várias folhas, mudando um pouco a posição dos braços em cada desenho. Ao folhear depressa, ela parece se mexer! Cada fantasia funciona como uma dessas imagens. A pausa controla quanto tempo você vê cada uma antes da troca.
\end{pensaassim}

\begin{codigo}[Receita da animação]
\begin{pilha}
\bk[scEvents]{quando \bandeira\ for clicado}
\sbc[scControl]{sempre}
\bk[scLooks]{próxima fantasia}
\bk[scControl]{espere \entrada{0.2} segundos}
\sbcend
\end{pilha}
\end{codigo}

\begin{dica}
Quanto \textbf{menor} a pausa, mais \textbf{rápida} a animação. Tente \entrada{0.05} para algo veloz e \entrada{0.5} para algo lento.
\end{dica}

\section{Andando e animando ao mesmo tempo}

Para o personagem andar de verdade, juntamos \textbf{movimento} e \textbf{fantasia}. A cada passo, o gatinho muda de perna.

\begin{pensaassim}[Pense assim: marchar parado ou sair do lugar]
Você pode marchar sem sair do lugar ou marchar avançando pelo chão. Trocar as fantasias pode mostrar as pernas se mexendo, mas não muda a posição do ator. Para ele atravessar o palco, combine a troca de fantasias com os blocos de movimento.
\end{pensaassim}

\begin{codigo}[Andar com animação]
\begin{pilha}
\bk[scEvents]{quando a tecla \menu{seta para a direita} for pressionada}
\bk[scMotion]{aponte para a direção \entrada{90}}
\bk[scMotion]{mova \entrada{10} passos}
\bk[scLooks]{próxima fantasia}
\end{pilha}
\end{codigo}

\begin{missao}[Missão 9: Gato andarilho]
Faça o gatinho andar para os dois lados, virando para o lado certo e trocando de perna a cada passo.
\end{missao}

\begin{codigo}[Esquerda e direita]
\begin{pilha}
\bk[scEvents]{quando a tecla \menu{seta direita} for pressionada}
\bk[scMotion]{aponte para a direção \entrada{90}}
\bk[scMotion]{mova \entrada{10} passos}
\bk[scLooks]{próxima fantasia}
\end{pilha}
\begin{pilha}
\bk[scEvents]{quando a tecla \menu{seta esquerda} for pressionada}
\bk[scMotion]{aponte para a direção \entrada{$-90$}}
\bk[scMotion]{mova \entrada{10} passos}
\bk[scLooks]{próxima fantasia}
\end{pilha}
\end{codigo}

\begin{cuidado}
Se o gatinho andar de ponta-cabeça ao virar para a esquerda, ajuste o \textbf{estilo de rotação} para \textbf{esquerda-direita} com o bloco \bloco[scMotion]{defina o estilo de rotação \menu{esquerda-direita}}.
\end{cuidado}

\section{Deslizando}

\bloco[scMotion]{deslize por \entrada{1} seg até x: \entrada{100} y: \entrada{0}} faz um movimento \textbf{suave} entre dois pontos, sem precisar de laço.

\begin{curiosidade}
Os animadores de filmes de desenho usam 12 a 24 desenhos \textbf{por segundo}. No Scratch o projeto roda a 30 quadros por segundo. Por isso a pausa de 0.1 segundo já parece um movimento natural!
\end{curiosidade}

\section{Hora de praticar}

\begin{exercicio}[Gato dançarino]
Faça o gatinho trocar de fantasia 10 vezes, com uma pausa de 0.2 segundo.
\end{exercicio}

\begin{solucao}
\begin{pilha}
\bk[scEvents]{quando \bandeira\ for clicado}
\sbc[scControl]{repita \entrada{10} vezes}
\bk[scLooks]{próxima fantasia}
\bk[scControl]{espere \entrada{0.2} segundos}
\sbcend
\end{pilha}
\end{solucao}

\begin{exercicio}[Vai e volta suave]
Faça o gatinho deslizar para a direita (x: 150) em 1 segundo e depois voltar ao centro (x: 0) em 1 segundo.
\end{exercicio}

\begin{solucao}
\begin{pilha}
\bk[scEvents]{quando \bandeira\ for clicado}
\bk[scMotion]{deslize por \entrada{1} seg até x: \entrada{150} y: \entrada{0}}
\bk[scMotion]{deslize por \entrada{1} seg até x: \entrada{0} y: \entrada{0}}
\end{pilha}
\end{solucao}

\begin{exercicio}[Cresce e encolhe]
Faça o gatinho crescer de 10 em 10 por 5 vezes e depois encolher de volta.
\end{exercicio}

\begin{solucao}
\begin{pilha}
\bk[scEvents]{quando \bandeira\ for clicado}
\sbc[scControl]{repita \entrada{5} vezes}
\bk[scLooks]{mude o tamanho em \entrada{10}}
\bk[scControl]{espere \entrada{0.1} segundos}
\sbcend
\sbc[scControl]{repita \entrada{5} vezes}
\bk[scLooks]{mude o tamanho em \entrada{$-10$}}
\bk[scControl]{espere \entrada{0.1} segundos}
\sbcend
\end{pilha}
\end{solucao}

\begin{desafio}[Passeio no palco]
Faça o gatinho andar para a direita com a seta direita, para a esquerda com a seta esquerda, e \textbf{pular} (subir 50 e voltar) com a seta para cima.
\end{desafio}

\begin{solucao}
\begin{pilha}
\bk[scEvents]{quando a tecla \menu{seta para cima} for pressionada}
\bk[scMotion]{mude y em \entrada{50}}
\bk[scControl]{espere \entrada{0.3} segundos}
\bk[scMotion]{mude y em \entrada{$-50$}}
\end{pilha}
As pilhas das setas esquerda e direita são as da Missão 9.
\end{solucao}
